\documentclass[11pt,a4paper]{article}

\usepackage[utf8]{inputenc}
\usepackage[T1]{fontenc}
\usepackage[french]{babel}
\usepackage[a4paper,margin=2.3cm]{geometry}
\usepackage{amsmath,amssymb,amsthm,mathtools}
\usepackage{lmodern}
\usepackage{enumitem}
\usepackage{hyperref}
\usepackage{xcolor}
\usepackage{fancyhdr}
\usepackage{stmaryrd}
\usepackage{mathrsfs}

\hypersetup{
    colorlinks=true,
    linkcolor=black,
    urlcolor=blue,
    citecolor=black
}

\setlength{\parindent}{0pt}
\setlength{\parskip}{0.6em}

\pagestyle{fancy}
\fancyhf{}
\fancyhead[L]{Algèbre linéaire --- Chapitre 2}
\fancyhead[R]{Pratique calculatoire}
\fancyfoot[C]{\thepage}

% Environnements
\newtheorem{definition}{Définition}[section]
\newtheorem{theoreme}[definition]{Théorème}
\newtheorem{proposition}[definition]{Proposition}
\newtheorem{corollaire}[definition]{Corollaire}
\newtheorem{lemme}[definition]{Lemme}
\theoremstyle{remark}
\newtheorem{remarque}[definition]{Remarque}
\newtheorem{exemple}[definition]{Exemple}
\newtheorem*{methode}{Méthode}
\newtheorem*{idee}{Idée}

% Commandes
\newcommand{\R}{\mathbb{R}}
\newcommand{\N}{\mathbb{N}}
\newcommand{\Z}{\mathbb{Z}}
\newcommand{\Q}{\mathbb{Q}}
\newcommand{\C}{\mathbb{C}}
\newcommand{\K}{\mathbb{K}}

\begin{document}

\begin{center}
    {\LARGE \textbf{Chapitre 2 --- Pratique calculatoire PGB}}\\[0.5em]
    {\large Outils algébriques fondamentaux pour l’algèbre linéaire}
\end{center}

\hrule
\vspace{1em}

\tableofcontents

\newpage

\section{Pourquoi ce chapitre est important ?}

En algèbre linéaire, on manipule sans cesse :
\begin{itemize}
    \item des expressions littérales,
    \item des égalités à transformer proprement,
    \item des sommes à simplifier,
    \item des produits à factoriser,
    \item des fractions à mettre au même dénominateur,
    \item des polynômes à développer ou à reconnaître.
\end{itemize}

Une grande partie des difficultés rencontrées dans les chapitres suivants ne vient pas des idées profondes, mais de calculs mal tenus :
\begin{itemize}
    \item oubli de parenthèses,
    \item erreur de signe,
    \item mauvaise factorisation,
    \item confusion entre une identité vraie pour tout \(x\) et une équation à résoudre,
    \item simplification illégitime par une quantité pouvant être nulle.
\end{itemize}

Ce chapitre a pour but de poser des bases solides de calcul.

\section{Égalités, identités et équations}

\subsection{Égalité et identité}

\begin{definition}
Une \textbf{égalité} est une affirmation de la forme
\[
A=B,
\]
où \(A\) et \(B\) désignent deux objets mathématiques.

Une \textbf{identité} est une égalité vraie pour toutes les valeurs des variables considérées dans un certain domaine.
\end{definition}

\begin{exemple}
L’égalité
\[
(x+y)^2=x^2+2xy+y^2
\]
est une identité : elle est vraie pour tous réels \(x\) et \(y\).

En revanche,
\[
x^2=4
\]
n’est pas une identité : elle n’est vraie que pour certains \(x\).
\end{exemple}

\begin{remarque}
En pratique, il faut toujours distinguer :
\begin{itemize}
    \item une formule valable \textbf{pour tout} \(x\),
    \item une équation que l’on cherche à résoudre.
\end{itemize}
\end{remarque}

\subsection{Équations}

\begin{definition}
Résoudre une \textbf{équation} consiste à déterminer toutes les valeurs de l’inconnue pour lesquelles l’égalité est vraie.
\end{definition}

\begin{exemple}
Résoudre
\[
x^2-1=0
\]
revient à déterminer les réels \(x\) tels que
\[
x^2=1,
\]
c’est-à-dire
\[
x=1 \quad \text{ou} \quad x=-1.
\]
\end{exemple}

\begin{methode}
Devant une équation, il faut toujours se demander :
\begin{itemize}
    \item dans quel ensemble on travaille ;
    \item si l’on peut factoriser ;
    \item si l’on peut faire apparaître une identité remarquable ;
    \item si certaines transformations risquent d’introduire ou de perdre des solutions.
\end{itemize}
\end{methode}

\section{Règles fondamentales de calcul}

\subsection{Calcul avec les parenthèses}

\begin{proposition}
Pour tous réels \(a,b,c\),
\[
a(b+c)=ab+ac.
\]
On parle de \textbf{distributivité} de la multiplication sur l’addition.
\end{proposition}

\begin{proof}
Cette propriété est fondamentale et découle des axiomes du calcul dans \(\R\). Elle est à la base de tous les développements et factorisations.
\end{proof}

\begin{corollaire}
Pour tous réels \(a,b,c\),
\[
(a+b)c=ac+bc,
\qquad
a(b-c)=ab-ac,
\qquad
(a-b)c=ac-bc.
\]
\end{corollaire}

\begin{exemple}
\[
3(x-2)=3x-6,
\qquad
-(a+b)=-a-b,
\qquad
-(a-b)=-a+b.
\]
\end{exemple}

\begin{remarque}
Le signe moins placé devant une parenthèse change tous les signes à l’intérieur :
\[
-(a+b)=-a-b,
\qquad
-(a-b)=-a+b.
\]
C’est une source d’erreurs très fréquente.
\end{remarque}

\subsection{Règles de simplification}

\begin{proposition}
Si \(a,b,c\in\R\), alors :
\[
a+c=b+c \Longrightarrow a=b.
\]
Si de plus \(c\neq 0\), alors
\[
ac=bc \Longrightarrow a=b.
\]
\end{proposition}

\begin{proof}
La première implication vient du fait qu’on peut soustraire \(c\) aux deux membres :
\[
a+c=b+c \Longrightarrow a+c-c=b+c-c \Longrightarrow a=b.
\]

La seconde vient du fait que si \(c\neq 0\), on peut diviser par \(c\) :
\[
ac=bc \Longrightarrow \frac{ac}{c}=\frac{bc}{c} \Longrightarrow a=b.
\]
\end{proof}

\begin{remarque}
On ne peut \textbf{jamais} simplifier par une quantité sans vérifier qu’elle est non nulle.
\end{remarque}

\begin{exemple}
L’égalité
\[
x(x-1)=x(x+2)
\]
ne permet pas de simplifier immédiatement par \(x\) sans précaution, car on peut avoir \(x=0\).

On écrit plutôt :
\[
x(x-1)-x(x+2)=0
\]
\[
x\big((x-1)-(x+2)\big)=0
\]
\[
x(-3)=0
\]
donc
\[
x=0.
\]
\end{exemple}

\section{Fractions}

\subsection{Définition et règles}

\begin{definition}
Si \(a\in\R\) et \(b\in\R^\ast=\R\setminus\{0\}\), la fraction
\[
\frac{a}{b}
\]
désigne le quotient de \(a\) par \(b\).
\end{definition}

\begin{proposition}
Soient \(a,b,c,d\in\R\) avec \(b\neq 0\) et \(d\neq 0\). Alors :
\[
\frac{a}{b}+\frac{c}{d}=\frac{ad+bc}{bd},
\]
\[
\frac{a}{b}-\frac{c}{d}=\frac{ad-bc}{bd},
\]
\[
\frac{a}{b}\cdot \frac{c}{d}=\frac{ac}{bd},
\]
\[
\frac{a}{b}\div \frac{c}{d}=\frac{a}{b}\cdot\frac{d}{c}=\frac{ad}{bc}
\qquad \text{si } c\neq 0.
\]
\end{proposition}

\begin{exemple}
\[
\frac{2}{3}+\frac{5}{4}=\frac{8}{12}+\frac{15}{12}=\frac{23}{12}.
\]
\end{exemple}

\subsection{Mise au même dénominateur}

\begin{methode}
Pour additionner ou soustraire des fractions, on commence par les écrire avec un dénominateur commun.
\end{methode}

\begin{exemple}
Simplifions
\[
\frac{1}{x}+\frac{1}{x+1}.
\]
Pour \(x\neq 0\) et \(x\neq -1\),
\[
\frac{1}{x}+\frac{1}{x+1}
=
\frac{x+1}{x(x+1)}+\frac{x}{x(x+1)}
=
\frac{2x+1}{x(x+1)}.
\]
\end{exemple}

\subsection{Fractions et simplification}

\begin{proposition}
Soient \(a,b,c\in\R\) avec \(b\neq 0\) et \(c\neq 0\). Alors
\[
\frac{ac}{bc}=\frac{a}{b}.
\]
\end{proposition}

\begin{proof}
Comme \(c\neq 0\), on peut diviser le numérateur et le dénominateur par \(c\), ce qui donne
\[
\frac{ac}{bc}=\frac{a}{b}.
\]
\end{proof}

\begin{remarque}
Là encore, la condition \(c\neq 0\) est essentielle.
\end{remarque}

\begin{exemple}
\[
\frac{x^2-1}{x-1}
=
\frac{(x-1)(x+1)}{x-1}
=
x+1
\qquad \text{si } x\neq 1.
\]
Cette simplification n’est valable que pour \(x\neq 1\).
\end{exemple}

\section{Puissances}

\subsection{Définitions et règles usuelles}

\begin{definition}
Soit \(a\in\R\) et \(n\in\N\). On définit
\[
a^0=1 \quad \text{si } a\neq 0,
\]
et pour \(n\geq 1\),
\[
a^n=\underbrace{a\times a\times \cdots \times a}_{n\ \text{facteurs}}.
\]
\end{definition}

\begin{proposition}
Pour \(a,b\in\R\) et \(m,n\in\N\),
\[
a^m a^n=a^{m+n},
\]
\[
(a^m)^n=a^{mn},
\]
\[
(ab)^n=a^n b^n.
\]
Si \(a\neq 0\), alors
\[
\frac{a^m}{a^n}=a^{m-n}
\qquad \text{si } m\geq n.
\]
\end{proposition}

\begin{proof}
Montrons par exemple
\[
a^m a^n=a^{m+n}.
\]
Le produit de \(m\) facteurs égaux à \(a\) par \(n\) autres facteurs égaux à \(a\) donne \(m+n\) facteurs égaux à \(a\), donc
\[
a^m a^n=a^{m+n}.
\]

Les autres propriétés se démontrent de façon analogue en comptant les facteurs.
\end{proof}

\begin{exemple}
\[
x^2x^5=x^7,
\qquad
(x^3)^4=x^{12},
\qquad
(2x)^3=8x^3.
\]
\end{exemple}

\subsection{Puissances négatives}

\begin{definition}
Si \(a\neq 0\) et \(n\in\N\), on définit
\[
a^{-n}=\frac{1}{a^n}.
\]
\end{definition}

\begin{exemple}
\[
2^{-3}=\frac{1}{2^3}=\frac{1}{8}.
\]
\end{exemple}

\begin{proposition}
Si \(a\neq 0\), alors pour tous entiers relatifs \(m,n\),
\[
a^m a^n=a^{m+n},
\qquad
(a^m)^n=a^{mn}.
\]
\end{proposition}

\begin{remarque}
Dès qu’il y a des puissances négatives, la condition \(a\neq 0\) devient indispensable.
\end{remarque}

\section{Valeur absolue}

\begin{definition}
Pour tout réel \(x\), la \textbf{valeur absolue} de \(x\), notée \(|x|\), est définie par
\[
|x|=
\begin{cases}
x & \text{si } x\geq 0,\\
-x & \text{si } x<0.
\end{cases}
\]
\end{definition}

\begin{remarque}
La valeur absolue mesure la distance de \(x\) à \(0\) sur la droite réelle.
\end{remarque}

\begin{proposition}
Pour tout \(x\in\R\),
\[
|x|\geq 0,
\qquad
|x|=0 \Longleftrightarrow x=0,
\qquad
|-x|=|x|,
\qquad
x^2=|x|^2.
\]
\end{proposition}

\begin{proof}
Les propriétés découlent immédiatement de la définition par cas.
\end{proof}

\begin{proposition}[Inégalité triangulaire]
Pour tous \(x,y\in\R\),
\[
|x+y|\leq |x|+|y|.
\]
\end{proposition}

\begin{proof}
On a
\[
(x+y)^2=x^2+2xy+y^2.
\]
Or
\[
2xy \leq 2|x||y|,
\]
donc
\[
(x+y)^2 \leq x^2+2|x||y|+y^2=(|x|+|y|)^2.
\]
Les deux membres étant positifs, on peut prendre les racines :
\[
|x+y|\leq |x|+|y|.
\]
\end{proof}

\begin{exemple}
\[
|3+(-5)|=|-2|=2
\quad \text{et} \quad
|3|+|-5|=8,
\]
ce qui vérifie bien
\[
|3+(-5)|\leq |3|+|-5|.
\]
\end{exemple}

\section{Identités remarquables}

\subsection{Formules de base}

\begin{proposition}[Identités remarquables]
Pour tous \(a,b\in\R\),
\[
(a+b)^2=a^2+2ab+b^2,
\]
\[
(a-b)^2=a^2-2ab+b^2,
\]
\[
(a+b)(a-b)=a^2-b^2.
\]
\end{proposition}

\begin{proof}
On développe :
\[
(a+b)^2=(a+b)(a+b)=a^2+ab+ba+b^2=a^2+2ab+b^2.
\]
De même,
\[
(a-b)^2=(a-b)(a-b)=a^2-ab-ba+b^2=a^2-2ab+b^2.
\]
Enfin,
\[
(a+b)(a-b)=a^2-ab+ab-b^2=a^2-b^2.
\]
\end{proof}

\begin{remarque}
Ces identités sont parmi les outils les plus utilisés de tout le programme.
\end{remarque}

\begin{exemple}
\[
x^2-9=x^2-3^2=(x-3)(x+3).
\]
\end{exemple}

\subsection{Applications classiques}

\begin{exemple}
Pour simplifier
\[
(x+1)^2-(x-1)^2,
\]
on développe :
\[
(x+1)^2=x^2+2x+1,
\qquad
(x-1)^2=x^2-2x+1.
\]
Donc
\[
(x+1)^2-(x-1)^2=(x^2+2x+1)-(x^2-2x+1)=4x.
\]
\end{exemple}

\begin{exemple}
Pour factoriser
\[
4x^2-25,
\]
on reconnaît une différence de deux carrés :
\[
4x^2-25=(2x)^2-5^2=(2x-5)(2x+5).
\]
\end{exemple}

\section{Développement et factorisation}

\subsection{Développer}

\begin{definition}
Développer une expression consiste à transformer un produit ou une puissance en somme de termes plus simples.
\end{definition}

\begin{exemple}
\[
(x+2)(x-3)=x^2-3x+2x-6=x^2-x-6.
\]
\end{exemple}

\subsection{Factoriser}

\begin{definition}
Factoriser une expression consiste à l’écrire sous la forme d’un produit.
\end{definition}

\begin{exemple}
\[
x^2-x=x(x-1).
\]
\end{exemple}

\begin{methode}
Pour factoriser, on pense surtout à :
\begin{itemize}
    \item mettre un facteur commun en évidence ;
    \item reconnaître une identité remarquable ;
    \item regrouper certains termes.
\end{itemize}
\end{methode}

\subsection{Mise en facteur commun}

\begin{proposition}
Pour tous \(a,b,c\in\R\),
\[
ab+ac=a(b+c).
\]
\end{proposition}

\begin{proof}
C’est exactement la distributivité écrite à l’envers.
\end{proof}

\begin{exemple}
\[
3x^2-6x=3x(x-2),
\]
\[
x^3+2x^2=x^2(x+2).
\]
\end{exemple}

\subsection{Factorisation par regroupement}

\begin{exemple}
Factorisons
\[
ax+ay+bx+by.
\]
On regroupe :
\[
ax+ay+bx+by=a(x+y)+b(x+y)=(a+b)(x+y).
\]
\end{exemple}

\section{Sommes usuelles}

\subsection{Somme des premiers entiers}

\begin{proposition}
Pour tout \(n\in\N\),
\[
1+2+\cdots+n=\frac{n(n+1)}{2}.
\]
\end{proposition}

\begin{proof}
La formule a déjà été démontrée au chapitre précédent par récurrence. Elle sera utilisée très souvent.
\end{proof}

\subsection{Somme des carrés}

\begin{proposition}
Pour tout \(n\in\N\),
\[
1^2+2^2+\cdots+n^2=\frac{n(n+1)(2n+1)}{6}.
\]
\end{proposition}

\begin{proof}
Cette formule peut se démontrer par récurrence. Elle est classique et très utile dans de nombreux calculs.
\end{proof}

\subsection{Somme géométrique}

\begin{proposition}
Soit \(q\in\R\) avec \(q\neq 1\). Pour tout \(n\in\N\),
\[
1+q+q^2+\cdots+q^n=\frac{1-q^{n+1}}{1-q}.
\]
\end{proposition}

\begin{proof}
Posons
\[
S=1+q+q^2+\cdots+q^n.
\]
Alors
\[
qS=q+q^2+\cdots+q^n+q^{n+1}.
\]
En soustrayant,
\[
S-qS=1-q^{n+1}.
\]
Donc
\[
(1-q)S=1-q^{n+1}.
\]
Comme \(q\neq 1\), on peut diviser par \(1-q\), d’où
\[
S=\frac{1-q^{n+1}}{1-q}.
\]
\end{proof}

\begin{exemple}
\[
1+2+2^2+2^3+2^4=\frac{1-2^5}{1-2}=31.
\]
\end{exemple}

\begin{remarque}
Si \(q=1\), alors
\[
1+q+\cdots+q^n=n+1.
\]
Il faut traiter ce cas à part, car la formule précédente ferait apparaître une division par \(0\).
\end{remarque}

\section{Binôme de Newton}

\subsection{Coefficients binomiaux}

\begin{definition}
Pour \(n\in\N\) et \(k\in\{0,\dots,n\}\), on définit le \textbf{coefficient binomial}
\[
\binom{n}{k}=\frac{n!}{k!(n-k)!}.
\]
\end{definition}

\begin{remarque}
Le nombre \(\binom{n}{k}\) compte le nombre de façons de choisir \(k\) objets parmi \(n\).
\end{remarque}

\subsection{Formule du binôme}

\begin{theoreme}[Binôme de Newton]
Pour tous \(a,b\in\R\) et tout \(n\in\N\),
\[
(a+b)^n=\sum_{k=0}^n \binom{n}{k}a^{n-k}b^k.
\]
\end{theoreme}

\begin{exemple}
Pour \(n=2\),
\[
(a+b)^2=\binom{2}{0}a^2+\binom{2}{1}ab+\binom{2}{2}b^2=a^2+2ab+b^2.
\]

Pour \(n=3\),
\[
(a+b)^3=a^3+3a^2b+3ab^2+b^3.
\]
\end{exemple}

\begin{remarque}
Dans la pratique courante en algèbre linéaire, on utilise surtout les cas \(n=2\) et \(n=3\), mais la formule générale devient utile dans l’étude des polynômes.
\end{remarque}

\section{Manipulation des expressions algébriques}

\subsection{Substitution et changement d’écriture}

\begin{exemple}
Si l’on pose
\[
u=x+1,
\]
alors
\[
(x+1)^2+3(x+1)+2=u^2+3u+2.
\]
Ce changement de variable temporaire peut simplifier les calculs.
\end{exemple}

\begin{remarque}
Savoir réécrire une expression sous une forme plus adaptée est souvent aussi important que savoir calculer.
\end{remarque}

\subsection{Réduction d’une expression}

\begin{definition}
Réduire une expression, c’est regrouper les termes de même nature.
\end{definition}

\begin{exemple}
\[
2x-3+5x+7=7x+4.
\]
\end{exemple}

\begin{remarque}
On ne peut regrouper que des termes comparables :
\[
2x+3x=5x,
\qquad
2x+3 \text{ ne se réduit pas.}
\]
\end{remarque}

\section{Premiers réflexes sur les polynômes}

\subsection{Définition informelle}

\begin{definition}
Une expression polynomiale en \(x\) est une somme finie de termes de la forme
\[
a_k x^k,
\]
où les \(a_k\) sont des coefficients et \(k\in\N\).
\end{definition}

\begin{exemple}
\[
P(x)=3x^4-2x+1
\]
est un polynôme.
\end{exemple}

\subsection{Degré}

\begin{definition}
Le \textbf{degré} d’un polynôme non nul est le plus grand exposant apparaissant avec un coefficient non nul.
\end{definition}

\begin{exemple}
Le polynôme
\[
P(x)=5x^3-2x+7
\]
est de degré \(3\).
\end{exemple}

\begin{remarque}
Le degré se lit sur la forme réduite du polynôme.
\end{remarque}

\subsection{Racines évidentes et factorisation}

\begin{proposition}
Si \(P(a)=0\), alors \(x-a\) est un facteur de \(P(x)\).
\end{proposition}

\begin{remarque}
Cette propriété sera démontrée rigoureusement dans le chapitre sur les polynômes. Pour l’instant, on l’utilise comme outil pratique sur des exemples simples.
\end{remarque}

\begin{exemple}
Considérons
\[
P(x)=x^2-5x+6.
\]
On vérifie :
\[
P(2)=4-10+6=0,
\qquad
P(3)=9-15+6=0.
\]
Ainsi \(2\) et \(3\) sont des racines, et
\[
P(x)=(x-2)(x-3).
\]
\end{exemple}

\section{Résolution d’équations algébriques simples}

\subsection{Équation produit nul}

\begin{proposition}
Pour tous réels \(a,b\),
\[
ab=0 \Longleftrightarrow a=0 \text{ ou } b=0.
\]
\end{proposition}

\begin{proof}
Si \(a=0\) ou \(b=0\), alors évidemment \(ab=0\).

Réciproquement, supposons \(ab=0\). Si \(a\neq 0\), alors on peut diviser par \(a\), ce qui donne
\[
b=0.
\]
Donc nécessairement \(a=0\) ou \(b=0\).
\end{proof}

\begin{exemple}
Résolvons
\[
(x-2)(x+5)=0.
\]
On obtient
\[
x-2=0 \quad \text{ou} \quad x+5=0,
\]
donc
\[
x=2 \quad \text{ou} \quad x=-5.
\]
\end{exemple}

\subsection{Équation du second degré sous forme factorisée}

\begin{exemple}
Résolvons
\[
x^2-7x+12=0.
\]
On cherche deux nombres dont la somme vaut \(7\) et le produit \(12\) :
\[
3 \quad \text{et} \quad 4.
\]
Donc
\[
x^2-7x+12=(x-3)(x-4),
\]
et l’équation devient
\[
(x-3)(x-4)=0.
\]
Ainsi
\[
x=3 \quad \text{ou} \quad x=4.
\]
\end{exemple}

\section{Méthodes concrètes de calcul}

\subsection{Bien poser un calcul}

\begin{methode}
Dans un calcul algébrique, il faut :
\begin{enumerate}
    \item écrire les étapes intermédiaires ;
    \item conserver les parenthèses tant qu’elles sont utiles ;
    \item annoncer la méthode utilisée : développer, factoriser, mettre au même dénominateur, etc. ;
    \item surveiller les conditions : dénominateur non nul, possibilité de simplification.
\end{enumerate}
\end{methode}

\subsection{Choisir entre développer et factoriser}

\begin{remarque}
Développer est utile quand :
\begin{itemize}
    \item on veut réduire l’expression ;
    \item on veut comparer des coefficients ;
    \item on veut additionner ou soustraire des expressions.
\end{itemize}

Factoriser est utile quand :
\begin{itemize}
    \item on veut résoudre une équation ;
    \item on veut faire apparaître un produit nul ;
    \item on veut simplifier une fraction ;
    \item on veut reconnaître une structure.
\end{itemize}
\end{remarque}

\begin{exemple}
Pour résoudre
\[
x^2-4=0,
\]
il est plus intelligent de factoriser :
\[
x^2-4=(x-2)(x+2).
\]
\end{exemple}

\section{Erreurs classiques à éviter}

\begin{itemize}
    \item Oublier les parenthèses :
    \[
    -(a-b)\neq -a-b \quad \text{mais} \quad -(a-b)=-a+b.
    \]

    \item Simplifier à tort :
    \[
    \frac{x^2+x}{x}=x+1
    \]
    n’est vrai que pour \(x\neq 0\).

    \item Croire que
    \[
    (a+b)^2=a^2+b^2.
    \]
    C’est faux en général ; il manque le terme \(2ab\).

    \item Confondre
    \[
    \frac{a+b}{c}
    \quad \text{et} \quad
    \frac{a}{c}+\frac{b}{c}
    \]
    qui est vrai,
    avec
    \[
    \frac{a}{b+c}
    \neq
    \frac{a}{b}+\frac{a}{c}.
    \]

    \item Diviser par une quantité pouvant être nulle.

    \item Perdre les conditions d’existence dans les calculs de fractions.
\end{itemize}

\section{Résumé du chapitre}

\begin{itemize}
    \item Une identité est une égalité vraie pour toutes les valeurs considérées.
    \item Une équation doit être résolue dans un ensemble donné.
    \item Les règles de distributivité sont à la base des développements et factorisations.
    \item Toute simplification par une quantité suppose qu’elle soit non nulle.
    \item Les fractions demandent une gestion rigoureuse des dénominateurs.
    \item Les règles sur les puissances doivent être parfaitement maîtrisées.
    \item Les identités remarquables permettent de développer vite et de factoriser intelligemment.
    \item Les formules de sommes usuelles servent constamment dans les calculs plus avancés.
    \item Savoir reconnaître une structure est souvent plus important que calculer brutalement.
\end{itemize}

\section*{Exercices d’application immédiate}

\begin{enumerate}[label=\textbf{Exercice \arabic*.}]
    \item Développer et réduire :
    \[
    (2x-3)(x+5), \qquad (x+1)^2-(x-2)^2.
    \]

    \item Factoriser :
    \[
    3x^2-12x, \qquad x^2-16, \qquad ax+ay+bx+by.
    \]

    \item Simplifier, en précisant les conditions d’existence :
    \[
    \frac{x^2-1}{x-1}, \qquad \frac{1}{x}+\frac{1}{x+2}.
    \]

    \item Résoudre dans \(\R\) :
    \[
    (x-4)(2x+3)=0,
    \qquad
    x^2-9=0.
    \]

    \item Montrer que pour tout \(n\in\N\),
    \[
    1+2+\cdots+n=\frac{n(n+1)}{2}.
    \]

    \item Montrer que pour tout réel \(q\neq 1\) et tout \(n\in\N\),
    \[
    1+q+q^2+\cdots+q^n=\frac{1-q^{n+1}}{1-q}.
    \]
\end{enumerate}

\end{document}